\section{Pointwise comparison through buffers}\label{fd:section}

We bound the probability of every local spin configuration after changing
prescribed spins, uniformly over all common pins.
We prove a finite-graph estimate in terms of a zero-field FK connection
probability, and then use critical crossings to control that probability
through lattice collars. All constants in the resulting geometric bounds
are independent of the number of lattice vertices.

\subsection{A finite-graph bound from extreme log odds}

We first allow any finite graph, nonnegative pair couplings, and arbitrary
finite fields. On a binary cube, a positive weight $p$ satisfies the
\emph{lattice condition}, or is \emph{log-supermodular}, if
\[
 p(x\wedge y)p(x\vee y)\ge p(x)p(y)
 \qquad\text{for every }x,y,
\]
where the minimum and maximum are taken coordinatewise.

The preservation of this condition under marginalization is a basic
property of multivariate total positivity \cite{KarlinRinott1980}.
We give the finite binary argument and then show that every log-odds
rectangle is bounded by the rectangle between the four constant terminal
configurations. This reduces the likelihood-ratio problem to two spins.

\begin{lemma}[Extreme rectangles in a ferromagnetic marginal]
\label{fd:mtp2}
Let \(p(a,b)>0\) be a marginal of a finite ferromagnetic Ising law with
arbitrary fields, on two disjoint sets of spins.  Set \(f=\log p\) and
\[
 L=f(+,+)+f(-,-)-f(+,-)-f(-,+).
\]
For every four configurations \(a,a',b,b'\),
\begin{equation}
 \left|f(a,b)+f(a',b')-f(a,b')-f(a',b)\right|\le L.
 \label{fd:extreme-rectangle}
\end{equation}
\end{lemma}

\begin{proof}
The Ising weight is log-supermodular: each ferromagnetic pair interaction
has nonnegative mixed difference and each field has zero mixed difference.
We recall why summing out spins preserves this property.  It suffices to
check a two-coordinate square after summing one other binary coordinate.
Denote the eight unsummed weights by
\(a,b,c,d,e,f,g,h\), where \((a,b,c,d)\) and \((e,f,g,h)\)
are the two square slices, in order \((--,+-,-+,++)\).
Log-supermodularity gives
\(ad\ge bc\), \(eh\ge fg\),
\(ah\ge bg\), and \(ah\ge cf\).
Moreover \((ah)(ed)=(ad)(eh)\ge(bg)(cf)\).
For positive numbers \(X,Y,U,V\), the conditions
\(X\ge\max(U,V)\) and \(XY\ge UV\) imply
\(X+Y\ge U+V\): assume \(U\ge V\) and use
\(X+UV/X\ge U+V\) for \(X\ge U\).
Applying this with \((X,Y,U,V)=(ah,ed,bg,cf)\) gives
\(ah+ed\ge bg+cf\).  Hence
\((a+e)(d+h)\ge(b+f)(c+g)\).
Nonnegative adjacent-coordinate square differences are equivalent to
log-supermodularity on a Boolean cube, by summing those differences along
coordinate paths.  Iterating the one-coordinate elimination proves that
the marginal \(p\) is log-supermodular.

For ordered \(u\le v\) and \(s\le t\), put
\[
 \Delta(u,v;s,t)=f(u,s)+f(v,t)-f(u,t)-f(v,s)\ge0.
\]
If \(D=f(a,b)+f(a',b')-f(a,b')-f(a',b)\), then the exact identity
\[
\begin{split}
L-D={}&\Delta(-,a';-,b)+\Delta(a,+;b',+)\\
      &+\Delta(a',+;-,b')+\Delta(-,a;b,+)
\end{split}
\]
proves \(D\le L\).  Interchanging \(b,b'\) proves \(-D\le L\).
No order between \(a,a'\) or between \(b,b'\) is required.
\end{proof}

The extreme rectangle has a direct Ising interpretation: restrict all
spins of each terminal set to a common sign, thereby contracting that set
to one vertex. The resulting two-spin interaction is unchanged by shifts
of the terminal fields. Balancing those fields will let us apply the
arbitrary-field covariance inequality of Ding--Song--Sun.

\begin{lemma}[Pointwise ratios with arbitrary common pins]
\label{fd:ratio}
Consider a finite ferromagnetic Ising graph, disjoint nonempty vertex sets
\(I,J\), and any fixed spins on a third set \(F\).
Delete \(F\), retaining its contribution as fields, and contract \(I\)
and \(J\) separately.  Let \(q_0\) be the zero-field FK probability that
the two contracted vertices are connected on this contracted, vertex-deleted
graph.  Then every two conditional spin laws on \(I\), obtained by
prescribing arbitrary configurations \(b,b'\) on \(J\), satisfy
\begin{equation}
 \exp(-4\operatorname{arctanh}q_0)
 \le\frac{\mu(\sigma_I=a\mid\sigma_J=b,\sigma_F)}
            {\mu(\sigma_I=a\mid\sigma_J=b',\sigma_F)}
 \le\exp(4\operatorname{arctanh}q_0)
 \quad\text{for every }a.
 \label{fd:pointwise-ratio}
\end{equation}
The same inequalities hold if either conditional law is replaced by any
mixture of these laws.  The assertion remains valid in the presence of
additional arbitrary finite external fields; those fields are set to zero
in the definition of \(q_0\).
\end{lemma}

\begin{proof}
A prescribed common spin $\tau_v$, $v\in F$, contributes
$K_{uv}\tau_v$ to the field at each surviving neighbor $u$.
After these fields have been included, remove the vertices of $F$.
Apply Lemma~\ref{fd:mtp2} to the joint marginal on \(I,J\).
Restricting both terminal sets to constant configurations is exactly their
separate contraction.  Its two-spin marginal has unnormalized weights
\[
 w(s,t)=\exp(C+Kst+As+Bt),\qquad s,t\in\{-1,+1\},
 \qquad K=L/4\ge0.
\]
The constants \(A,B\) include all effective fields from the other vertices.
Add external fields \(-A,-B\) at the two contracted vertices.  This changes
their marginal to weights proportional to \(e^{Kst}\), without changing
\(K\).  Both means now vanish and the covariance is \(\tanh K\).
The Ding--Song--Sun inequality applies to this finite ferromagnetic graph
with these arbitrary mixed-sign fields
\cite[Corollary~1.3]{DSS2023}.  It gives
\[
 \tanh K\le\langle\sigma_I\sigma_J\rangle_0=q_0.
\]
The last equality is the zero-field Edwards--Sokal identity
\cite{EdwardsSokal1988}.
Internal edges of a contracted terminal contribute constants and parallel
edges have nonnegative summed coupling, so all hypotheses of that
inequality are preserved.  Thus \(L\le4\operatorname{arctanh}q_0\).

For fixed \(b,b'\), Lemma~\ref{fd:mtp2} bounds the oscillation in \(a\)
of the logarithm of the conditional density ratio by \(L\).
The ratio has mean one under its denominator law; its minimum is therefore
at most one and its maximum at least one.  The ratio lies throughout
\([e^{-L},e^L]\), as claimed.  Averaging a pointwise inequality proves the
mixture assertion, first for one law and then for the other.
\end{proof}

\begin{remark}[Symmetric terminals]\label{fd:symmetric-terminals}
The extreme log-odds rectangle and the connection probability $q_0$
are unchanged when $I$ and $J$ are interchanged. Thus the pointwise
comparison applies with the observed and prescribed terminal sets
reversed. Common pins are still deleted from the comparison graph;
oppositely signed common pins are never wired together. This symmetry
will also be used when a small region of changed pins is the source
and the observed set is much larger.
\end{remark}

\subsection{Critical collars, small patches, and local charts}

We now set every retained nearest-neighbor coupling to $\beta_c$.
To bound the connection probability in Lemma~\ref{fd:ratio}, it suffices
to find a dual circuit separating the terminals. A fixed collar provides
a probability bounded away from zero; many disjoint collars give a power
of the ratio of scales. Common prescribed spins still contribute fields
in the Ising law and are deleted only in the zero-field comparator.

We use the random-cluster representation with
\(p_c=1-e^{-2\beta_c}\) and cluster weight \(2\).  An FK boundary wiring is
a partition of the boundary vertices.  An open crossing here always uses
actual lattice edges; an abstract boundary identification is not itself a
crossing.  Free clusters receive independent fair spin signs in the
Edwards--Sokal representation \cite{EdwardsSokal1988}, and a cluster attached to plus pins has sign
plus.
We use the FK domain Markov property, association, and monotonicity
under edge addition and boundary wiring in this finite-graph setting;
these properties are recalled in \cite[Section~2]{DHN2011}. Their
association and ordering principles originate in
\cite{FKG1971,Holley1974}.
Association applies to decreasing events too.  Thus overlapping
crossings used to construct a dual circuit can be combined by
association, while disjoint tests can also be combined conditionally.

\begin{lemma}[Conditional crossing tests]\label{fd:conditional-rsw}
Fix \(A\ge1\).  There is \(c_A>0\) such that a lattice rectangle with both
side lengths at least one and aspect ratio in \([A^{-1},A]\) has probability
at least \(c_A\) of either specified open crossing, and of either specified
dual crossing, at criticality under any boundary wiring.  Changing each
side by a bounded number of lattice spacings only changes the constant.

These lower bounds also apply to a decreasing dual-crossing test in a
vertex-deleted graph, conditional on all bonds outside its full rectangular
test region, provided no terminal identification has vertices in the
interior of that region.  Additional vertex deletions preserve the bound.
For finitely many tests with disjoint sets of unexposed edges, the conditional
bounds may be multiplied.
\end{lemma}

\begin{proof}
The rectangle assertion is the critical FK-Ising RSW theorem
\cite[arXiv version~1, Theorem~1]{DHN2011}, with planar duality for dual crossings.
The domain Markov property turns an exterior bond configuration into a
boundary wiring.  Restore the deleted vertices of the test rectangle,
restore their edges, and wire its entire outer boundary.  The resulting
measure stochastically dominates the conditional original bond law, so a
decreasing dual event is at least as likely in the original law.  The
domination follows, for example, from the one-edge conditional probabilities
\(p_c\) and \(p_c/(p_c+2(1-p_c))\), according as the endpoints are or are
not already connected: adding edges or wiring increases these probabilities.
Restoring isolated vertices first and then their edges has the same effect.
The hypothesis about terminals ensures that no identification in the
interior of the restored rectangle is being discarded.

For the last assertion, order the test regions.  Conditioning on all edges
outside the next region includes the results of all earlier tests.  Applying
the uniform bound and then averaging proves the product lower bound by
induction.  The corresponding product upper bound for simultaneous failure
of tests follows in the same way.  If a bounded lattice correction makes
a side have bounded length, finitely many prescribed closed or open edges,
using the same one-edge probabilities, supplies a uniform replacement test.
\end{proof}

\begin{corollary}[Buffered and small-patch comparisons]
\label{fd:buffered-ratio}\label{fd:patch-ratio}
Fix \(c>0\).  There is \(K_c<\infty\) such that, on the square lattice,
\(I\subseteq z+B_R\) and \(J\subseteq(z+B_{(1+c)R})^c\) imply
pointwise density ratios in \([K_c^{-1},K_c]\) for the conditional laws
in Lemma~\ref{fd:ratio}.  Other common pins may occur anywhere off
\(I\cup J\).

There are absolute \(K<\infty\) and \(\delta>0\) such that, if
\(I\subseteq z+B_w\), \(J\subseteq(z+B_R)^c\), and \(R\ge8w\ge8\),
the logarithm of every such conditional density ratio has absolute value
at most
\begin{equation}
 K(w/R)^\delta.
 \label{fd:patch-power}
\end{equation}
In particular the total variation distance is at most
\(K'(w/R)^\delta\), and the expectation difference of a bounded test
function is at most \(2K'\|h\|_\infty(w/R)^\delta\).
No free neighborhood of the source patch \(I\) is required.

The buffered comparison also holds between laws from different domains,
including plane, torus, or cylinder laws, whenever their common collar
is an isometric square-lattice chart and all differing data lie beyond it.
It holds after applying the same conditional randomization kernel to local
spins, in particular internal Edwards--Sokal bond sampling.
\end{corollary}

\begin{proof}
An annulus of fixed radius ratio separating the two terminals contains a
dual circuit with probability at least \(c_1(c)>0\).  Cover a middle
subannulus by a fixed number of overlapping rectangles, require dual
crossings in both directions of their overlap squares, and use
Lemma~\ref{fd:conditional-rsw}.  The crossings join to an actual dual
circuit, which separates the terminals even after their separate internal
wiring.  Extra common pins are deleted in the FK comparator; restoring
them only increases connectivity.  Thus \(q_0\le1-c_1(c)\), giving the
first assertion by Lemma~\ref{fd:ratio}.

Between \(B_{2w}\) and \(B_{R/2}\), choose annuli with radius ratio, say,
four and disjoint test edges.  There are at least
\(c_2\log(R/w)-c_3\) such annuli.  A connection must avoid the dual
circuit in each, so Lemma~\ref{fd:conditional-rsw} gives
\(q_0\le C_1(w/R)^{\delta_1}\).
The fixed annulus between the terminals also gives \(q_0\le1-c_4\).
On \([0,1-c_4]\), \(\operatorname{arctanh}q\le C_2q\).
This proves the logarithmic bound.  The total variation and expectation
bounds follow by integrating the density difference.

For different domains, condition on the exterior neighbor spins of a
square contained in the common chart.  Each local law is a mixture of
the same finite-volume conditional laws, so the mixture assertion applies.
Finite-volume approximation gives the plane and cylinder statements;
the bounds are uniform in the approximating volume.  Finally, applying
the same Markov kernel preserves pointwise domination of measures.
\end{proof}

\begin{corollary}[Local charts and finite patch covers]
\label{fd:finite-cover}
The preceding buffered and patch bounds remain valid in a chart clipped
by a free or a fixed physical boundary, provided the local graph and
the common prescribed boundary spins agree in the compared laws.
Deleted vertices and missing edges may occur in the chart. For laws on
different domains, assume that conditioning on spins beyond the common
chart gives the same local interaction and pinning rule.

More generally, suppose the observed set is partitioned into at most
$m$ patches. If each patch has a square-lattice chart separating it
from every changed datum by a fixed positive relative distance, the
whole observed law has a pointwise density comparison with constant
equal to the product of the patch constants. The same conclusion holds
if instead the changed source set has such a partition, with each
source patch separated from the whole observed set. Constants depend
on $m$ and the fixed chart clearances. Polygonal and rotated collars
are therefore allowed whenever they admit such a finite cover.
\end{corollary}

\begin{proof}
First fix a common finite local conditional graph. In its zero-field
FK comparator restore any vertices and edges missing from the full
square test rectangles. Restoration increases connectivity. The
decreasing dual tests in Lemma~\ref{fd:conditional-rsw} therefore give
the same upper bound on $q_0$ as before. Prescribed physical-boundary
vertices are among the common pins: delete them in the comparator
before this restoration. No identification is added in a test
rectangle's interior. Lemma~\ref{fd:ratio} proves the asserted
pointwise bounds for this conditional graph.

For different volumes or continuations, condition at the edge of the
common chart. Both laws are mixtures of the same local kernels, with
possibly different mixing measures. Apply the mixture assertion of
Lemma~\ref{fd:ratio}. The restoration argument bounds the auxiliary
connection probability; the common conditional kernel is what permits
the comparison between volumes.

For a cover of the observed set, reveal its patches in a fixed order
and factor each configuration probability by the chain rule. In the
comparison for the next patch, previously revealed patches are common
pins. They can lie anywhere off the two terminals. Each conditional
ratio has its stated bound, and multiplication gives the product
constant. For a cover of the source, telescope the changes in its
patches; all other source patches are common pins. Apply
Remark~\ref{fd:symmetric-terminals} to place the small patch in the
inner terminal. Mixtures over source assignments and common
randomizations preserve the resulting bounds.
\end{proof}
